\htmlhr
\chapterAndLabel{PICO Checker}{pico-checker}

The PICO Checker enforces practical immutability for classes and objects.
It distinguishes references that permit mutation from references that are
read-only, and it supports classes whose instances may be used either
mutably or immutably.
The PICO type system is described by Ta~\cite{Ta2018},
Sun~\cite{Sun2021}, and
\href{https://uwspace.uwaterloo.ca/items/990ef07b-60be-420c-98ac-321ee7ad071c}{Shi}.

PICO is useful for APIs that expose read-only references to objects, for classes
that should not mutate their abstract state after construction, and for
code that wants to express when a method may mutate its receiver.

To run the PICO Checker, supply the
\code{-processor org.checkerframework.checker.pico.PICOChecker}
command-line option to javac.

\sectionAndLabel{PICO qualifiers}{pico-qualifiers}

The main PICO type qualifiers are:

\begin{figure}
\includeimage{pico}{5cm}
\caption{Type hierarchy for the PICO mutability type system.
  The \refqualclass{checker/pico/qual}{PolyMutable} qualifier is omitted
  because it is substituted with a concrete qualifier at each use.
  The internal \refqualclass{checker/pico/qual}{PICOLost} qualifier is
  omitted.}
\label{fig-pico-hierarchy}
\end{figure}

\begin{description}

\item[\refqualclass{checker/pico/qual}{Readonly}]
  is the top qualifier in the PICO hierarchy.
  A reference of type \<@Readonly T> does not permit mutation of the
  referenced object's fields through that reference.
  The object might still be mutated through another alias with a more
  specific mutability type.

\item[\refqualclass{checker/pico/qual}{Mutable}]
  indicates that a reference permits mutation of the referenced object's
  fields.  When written on a class declaration, it means that instances of
  the class have a mutable declaration bound.

\item[\refqualclass{checker/pico/qual}{Immutable}]
  indicates that the referenced object is immutable.  When written on a
  class declaration, it means that instances of the class have an immutable
  declaration bound.

\item[\refqualclass{checker/pico/qual}{ReceiverDependentMutable}]
  indicates that the effective mutability depends on the receiver.
  On a member type, it is replaced by the mutability of the receiver at a
  use site.  On a class declaration, it means that the class may be used at
  different mutability instantiations, such as \<@Mutable> or \<@Immutable>.

\item[\refqualclass{checker/pico/qual}{PolyMutable}]
  is the polymorphic qualifier for the PICO mutability hierarchy.
  At a call site, occurrences of \<@PolyMutable> in one polymorphic use are
  substituted with the same concrete mutability qualifier.

\end{description}

The checker also uses
\refqualclass{checker/pico/qual}{PICOLost} internally when viewpoint
adaptation loses precise mutability information, and
\refqualclass{checker/pico/qual}{PICOBottom} as the bottom qualifier.
Programmers should normally write one of the main qualifiers above instead.

\sectionAndLabel{Class bounds and receiver-dependent mutability}{pico-class-bounds}

A PICO qualifier on a class declaration is a bound on valid uses of that
class.  For example, a class declared \<@Immutable> may only be used through
immutable-compatible references, while a class declared
\<@ReceiverDependentMutable> may be used at multiple mutability
instantiations.

\begin{Verbatim}
import org.checkerframework.checker.pico.qual.*;

@ReceiverDependentMutable class Buffer {
    private String contents;

    String get(@Readonly Buffer this) {
        return contents;
    }

    void set(@Mutable Buffer this, String contents) {
        this.contents = contents;
    }
}
\end{Verbatim}

The receiver annotations on \<get> and \<set> describe whether the method
may mutate the receiver.  A \<@Readonly> receiver can call \<get>, but it
cannot call \<set>.

\begin{Verbatim}
void use(@Mutable Buffer mutable, @Readonly Buffer readonly) {
    mutable.set("new value");
    mutable.get();

    readonly.get();
    readonly.set("new value"); // error: method.invocation.invalid
}
\end{Verbatim}

\sectionAndLabel{Object creation}{pico-object-creation}

The mutability of a newly-created receiver-dependent object can be selected
by the target type or by an explicit annotation on the creation expression.

\begin{Verbatim}
@Mutable Buffer mutable = new @Mutable Buffer();
@Immutable Buffer immutable = new @Immutable Buffer();
\end{Verbatim}

Creating a \<@Readonly> object is not permitted, because \<@Readonly>
describes a reference that cannot be used for mutation; it does not describe
a concrete object mutability.

\sectionAndLabel{Assignable fields}{pico-assignable-fields}

PICO distinguishes an object's abstract state from fields that may be
reassigned after initialization.  A field annotated
\refqualclass{checker/pico/qual}{Assignable} may be reassigned even when
ordinary fields of the receiver may not be modified.

\begin{Verbatim}
import org.checkerframework.checker.pico.qual.*;

@ReceiverDependentMutable class CachedName {
    private final String name;
    @Assignable private int cachedHash;

    CachedName(String name) {
        this.name = name;
    }

    int cachedHash(@Readonly CachedName this) {
        if (cachedHash == 0) {
            cachedHash = name.hashCode();
        }
        return cachedHash;
    }
}
\end{Verbatim}

Use \<@Assignable> for cache fields or other representation details that
are not part of the object's abstract state.

\sectionAndLabel{Implicitly immutable types}{pico-implicitly-immutable-types}

PICO treats primitive values, string literals, and common wrapper/value
classes such as \<String>, \<Integer>, and \<BigInteger> as immutable.
Such types should be used with \<@Immutable> or with the default qualifier
that PICO assigns to them.

% LocalWords:  PICO mutably mutability PolyMutable PICOLost PICOBottom
% LocalWords:  BigInteger
